\subsection{总体结果与对比分析}
\label{sec:overall_results_excerpt}

表 \ref{tab:main_results} 和图 \ref{fig:main_results} 汇报了不同规模下的总体结果。修正后的实验结果支持三个更稳健的结论。第一，在可处理实例上，CG 与精确基准保持一致：当 $N=20$ 时，CG、MILP60 与 MILP300 的平均目标值均为 24666.41；当 $N=50$ 时，CG 与 MILP300 的平均目标值均为 82695.01，而 MILP60 在 60s 时限下只有 80\% 的实例成功求解。第二，当规模增加到 $N=100$ 时，直接 MILP 已呈现出明显的不稳定性：MILP300 的成功率仅为 20\%，平均运行时间达到 290.1s，而 CG 仍以 100\% 成功率在 18.5s 内稳定求解。第三，在 $N=200$ 和 $N=500$ 下，MILP 已无法提供有效基准，而 CG 仍分别在 27.9s 和 72.5s 内完成求解，显示出良好的可扩展性。

相较于启发式规则，CG 在所有规模下均能取得更低成本。以 Greedy 为参照，CG 在 $N=20,50,100,200,500$ 下的平均降本幅度分别约为 9.7\%、2.9\%、16.6\%、1.3\% 和 2.5\%；相对于 FIFO，降本幅度分别约为 9.8\%、2.8\%、16.4\%、1.0\% 和 2.7\%。其中 $N=100$ 时改进最为显著，说明在中等拥挤水平下，前瞻性的资源协调最能体现价值。
